%% SwimAth — Diplomová práce
%% Návrh a implementace webového systému pro automatickou analýzu plavecké techniky z videa
%% Autor: Bc. Vít Vagner
%% Kompilace: xelatex main.tex && bibtex main && xelatex main.tex && xelatex main.tex

%%%%%%%%%%%%%%%%%%%%%%%%%%%%%%%%%%%%%%%%%
% CLASS OPTIONS
% language: czech/english/slovak
% thesis type: bachelor/master/dissertation
% electronic (oneside) or printed (twoside)
%%%%%%%%%%%%%%%%%%%%%%%%%%%%%%%%%%%%%%%%%
\documentclass[czech,master,twoside]{ctufit-thesis}

%%%%%%%%%%%%%%%%%%%%%%%%%%%%%%%%%%
% FILL IN THIS INFORMATION
%%%%%%%%%%%%%%%%%%%%%%%%%%%%%%%%%%
\ctufittitle{Návrh a~implementace webového systému pro automatickou analýzu plavecké techniky z~videa}
\ctufitauthorfull{Bc.~Vít Vagner}
\ctufitauthorsurnames{Vagner}
\ctufitauthorgivennames{Vít}
\ctufitsupervisor{TODO}  % TODO: doplnit vedoucího (musí mít Ph.D.)
\ctufitdepartment{Katedra softwarového inženýrství}  % TODO: ověřit katedru
\ctufitprogram{Informatika}
\ctufitspecialisation{Softwarové inženýrství}
\ctufityear{2027}
\ctufitdeclarationplace{Praze}
\ctufitdeclarationdate{\today}

\ctufitabstractCZE{%
Tato diplomová práce představuje SwimAth, webový systém pro automatickou analýzu plavecké techniky z~videa.
Práce se zaměřuje na~návrh a~implementaci kompletní softwarové architektury zahrnující asynchronní zpracování videa,
integraci ML pipeline pro odhad lidské pózy a~biomechanickou analýzu.
Systém zpracovává nahraný videozáznam plavce, extrahuje klíčové body těla pomocí fine-tuned modelu
(ViTPose/RTMPose), filtruje je přes ragdoll constraints, vypočítává biomechanické metriky
(úhly kloubů, rotace trupu, frekvence a~délka záběru) a~porovnává je s~referenčními vzory
pro dvě úrovně plavců pomocí Dynamic Time Warping.
Výstupem je skóre techniky, detekce konkrétních chyb a~personalizovaná doporučení ke~zlepšení.
Webová aplikace je postavena na~architektuře Vue.js~3 + FastAPI + Celery + PostgreSQL
s~důrazem na~škálovatelnost, testovatelnost a~bezpečnost (GDPR).
Práce zahrnuje návrh REST API, datového modelu, asynchronního zpracování úloh,
benchmark pose estimation modelů na~datasetu SwimXYZ a~vyhodnocení systému na~reálných videích.%
}

\ctufitabstractENG{%
This master thesis presents SwimAth, a~web-based system for automated swimming technique analysis from video.
The thesis focuses on~the design and implementation of~a~complete software architecture
encompassing asynchronous video processing, ML pipeline integration for human pose estimation,
and biomechanical analysis.
The system processes uploaded swimmer video recordings, extracts body keypoints using a~fine-tuned model
(ViTPose/RTMPose), filters them through ragdoll constraints, computes biomechanical metrics
(joint angles, trunk rotation, stroke rate and length), and compares them with reference patterns
for two skill levels using Dynamic Time Warping.
The output includes a~technique score, detection of~specific errors, and personalized improvement recommendations.
The web application is built on~a~Vue.js~3 + FastAPI + Celery + PostgreSQL architecture
with emphasis on~scalability, testability, and security (GDPR).
The thesis includes REST API design, data model, asynchronous task processing,
a~benchmark of~pose estimation models on~the SwimXYZ dataset, and system evaluation on~real-world videos.%
}

\ctufitkeywordsCZE{softwarová architektura, webová aplikace, analýza plavecké techniky, odhad lidské pózy, asynchronní zpracování, biomechanika plavání, Dynamic Time Warping, FastAPI, Vue.js, strojové učení}
\ctufitkeywordsENG{software architecture, web application, swimming technique analysis, human pose estimation, asynchronous processing, swimming biomechanics, Dynamic Time Warping, FastAPI, Vue.js, machine learning}
%%%%%%%%%%%%%%%%%%%%%%%%%%%%%%%%%%
% END FILL IN
%%%%%%%%%%%%%%%%%%%%%%%%%%%%%%%%%%

%%%%%%%%%%%%%%%%%%%%%%%%%%%%%%%%%%
% CUSTOMIZATION of this template
%%%%%%%%%%%%%%%%%%%%%%%%%%%%%%%%%%

\RequirePackage{iftex}[2020/03/06]
\iftutex % XeLaTeX and LuaLaTeX
    \RequirePackage{ellipsis}[2020/05/22]
\else
    \errmessage{Only compilation with XeLaTeX or LuaLaTeX is allowed}
    \stop
\fi

% hyperlinks
\hypersetup{
    pdfpagelayout=TwoPageRight,
    colorlinks=true,
    linkcolor=black,
    citecolor=blue!60!black,
    urlcolor=blue!70!black,
    pdfborder={0 0 0.1}
}

\RequirePackage{pdfpages}[2020/01/28]

%%%%%%%%%%%%%%%%%%%%%%%%%%%%%%%%%%
% CUSTOMIZATION END
%%%%%%%%%%%%%%%%%%%%%%%%%%%%%%%%%%


%%%%%%%%%%%%%%%%%%%%%%
% PACKAGES
%%%%%%%%%%%%%%%%%%%%%%

% --- Bibliography (biblatex with natbib compatibility) ---
\usepackage[backend=bibtex,style=numeric,natbib=true,sorting=none]{biblatex}
\addbibresource{references.bib}

% --- Graphics & Floats ---
\usepackage{graphicx}
\graphicspath{{figures/}}
\usepackage{float}
\usepackage{subcaption}

% --- Tables ---
\usepackage{booktabs}
\usepackage{tabularx}
\usepackage{longtable}

% --- TikZ & pgfplots ---
\usepackage{tikz}
\usepackage{pgfplots}
\pgfplotsset{compat=1.18}
\usetikzlibrary{patterns, positioning, arrows.meta}

% --- Code Listings ---
\usepackage{listings}
\lstset{
    basicstyle=\ttfamily\small,
    breaklines=true,
    frame=single,
    numbers=left,
    numberstyle=\tiny\color[gray]{0.5},
    tabsize=4,
    captionpos=b
}

% --- Units ---
\usepackage{siunitx}

% --- Misc ---
\usepackage{enumitem}
\usepackage{csquotes}
\providecommand{\uv}[1]{\enquote{#1}}

%%%%%%%%%%%%%%%%%%%%%%
% PACKAGES END
%%%%%%%%%%%%%%%%%%%%%%

\begin{document}
\frontmatter\frontmatterinit

\thispagestyle{empty}\maketitle\thispagestyle{empty}\cleardoublepage

% TODO: nahradit za PDF zadání z ProjectsFIT
% \includepdf[pages={1-}]{assignment.pdf}

\imprintpage
\stopTOCentries

%%%%%%%%%%%%%%%%%%%
% ACKNOWLEDGMENT
%%%%%%%%%%%%%%%%%%%
\begin{acknowledgmentpage}
	% TODO: Doplnit poděkování.
	Chtěl bych poděkovat \ldots
\end{acknowledgmentpage}

%%%%%%%%%%%%%%%%%%%
% DECLARATION
%%%%%%%%%%%%%%%%%%%
\begin{declarationpage}
	% TODO: Vybrat schválený text prohlášení z https://courses.fit.cvut.cz/SZZ/dokumenty/index.html
	Prohlašuji, že jsem předloženou práci vypracoval samostatně a~že jsem uvedl veškeré použité informační zdroje v~souladu s~Metodickým pokynem o~dodržování etických principů při přípravě vysokoškolských závěrečných prací.
\end{declarationpage}

\printtwopageabstract

%%%%%%%%%%%%%%%%%%%
% ABBREVIATIONS
%%%%%%%%%%%%%%%%%%%
\chapter{\thectufitabbreviationlabel}

\begin{tabular}{rl}
	API   & Application Programming Interface \\
	CI/CD & Continuous Integration / Continuous Deployment \\
	COCO  & Common Objects in Context \\
	CORS  & Cross-Origin Resource Sharing \\
	CV    & Computer Vision -- počítačové vidění \\
	DBA   & DTW Barycenter Averaging \\
	DTW   & Dynamic Time Warping \\
	FPS   & Frames Per Second \\
	GDPR  & General Data Protection Regulation \\
	GPU   & Graphics Processing Unit \\
	IdC   & Index of Coordination \\
	LSTM  & Long Short-Term Memory \\
	ML    & Machine Learning -- strojové učení \\
	NFR   & Non-Functional Requirement \\
	ONNX  & Open Neural Network Exchange \\
	ORM   & Object-Relational Mapping \\
	PE    & Pose Estimation -- odhad pózy \\
	REST  & Representational State Transfer \\
	SL    & Stroke Length -- délka záběru \\
	SR    & Stroke Rate -- frekvence záběru \\
	WS    & WebSocket \\
\end{tabular}

\resumeTOCentries

\tableofcontents
\listoffigures
\begingroup
\let\clearpage\relax
\listoftables
\endgroup

\mainmatter\mainmatterinit

%%%%%%%%%%%%%%%%%%%
% THE THESIS
%%%%%%%%%%%%%%%%%%%

%---------------------------------------------------------------
\chapter{Úvod}
\label{ch:uvod}
%---------------------------------------------------------------

% TODO: Rozepsat úvod — motivace, cíl práce, přínos, struktura textu
% Osnova:
% - Plavání jako sport s vysokými nároky na techniku
% - Dostupnost technologické analýzy (drahé systémy, málo dostupných nástrojů)
% - Cíl: end-to-end systém video → metriky → detekce chyb → zpětná vazba
% - Přínos: první veřejně dostupný pipeline kombinující PE + biomechaniku + personalizovaný feedback
% - Struktura práce (přehled kapitol)

%---------------------------------------------------------------
\chapter{Rešerše}
\label{ch:reserse}
%---------------------------------------------------------------

\section{Biomechanika plavání}
\label{sec:biomechanika}
% TODO: Styly, metriky, co dělá dobrého plavce
% Zdroje: Maglischo 2003, Craig & Pendergast 1979, Barbosa 2011
% - 4 plavecké styly (kraul, znak, prsa, motýlek)
% - Klíčové metriky: SR, SL, V = SR × SL, IdC, úhly kloubů, rotace trupu
% - Fáze záběru: catch, pull, push, recovery
% - Nejčastější chyby: dropped elbow (61.3% dle Barden & Kell 2014), špatná rotace trupu

\section{Odhad lidské pózy}
\label{sec:pose_estimation}
% TODO: Architektury: OpenPose → HRNet → ViTPose → RTMPose
% Zdroje: Xu 2022 (ViTPose), Jiang 2023 (RTMPose), Einfalt 2018/2019
% - Evoluce: bottom-up vs top-down
% - ViTPose: Vision Transformer, škálovatelný 20M-1B params
% - RTMPose: 75.8% AP COCO, 90+ FPS CPU, 430+ FPS GPU
% - Plavecky-specifické: Einfalt (Augsburg), domain-specific challenges

\section{Datasety pro plaveckou analýzu}
\label{sec:datasety}
% TODO: SwimXYZ, Augsburg, SwimmerNET, CADDY
% Zdroje: Fiche 2023 (SwimXYZ), Giulietti 2023 (SwimmerNET)
% - SwimXYZ: syntetický, 11,520 videí, 3.4M snímků, 48 keypointů
% - Augsburg: reálný, 4 styly, na žádost
% - SwimmerNET: podvodní, kraul, 2,021 snímků
% - Tabulka srovnání datasetů

\section{Strojové učení ve sportovní analýze}
\label{sec:ml_sport}
% TODO: Existující systémy, co chybí
% Diskuse RAG/LLM přístupu: Comendant 2024, Talking Tennis, BoxingPro
% Argumentace proč ML pipeline místo LLM:
% - Závislost na ext. API, nedeterminismus, latence, náklady, offline, reprodukovatelnost
% - Porovnání trade-offs: pravidlový systém vs. LLM-generated feedback

\section{Podvodní výzvy}
\label{sec:podvodni_vyzvy}
% TODO: Refrakce, bubliny, viditelnost
% - Optické artefakty pod vodou
% - Problémy s detekcí pózy (okluze, reflexy)
% - Existující přístupy: SwimmerNET, DeepLabCut

%---------------------------------------------------------------
\chapter{Návrh řešení}
\label{ch:navrh}
%---------------------------------------------------------------

\section{Architektura systému}
\label{sec:architektura}
% TODO: Vue.js 3 + FastAPI + Celery + Redis + PostgreSQL
% - Diagram architektury (figure)
% - Upload flow: video + úroveň + pohled + styl → WebSocket progress → výsledky
% - Asynchronní zpracování (Celery workers)

\section{Výběr modelu pro odhad pózy}
\label{sec:vyber_modelu}
% TODO: Porovnání MediaPipe vs RTMPose vs ViTPose
% - Důvody výběru (přesnost, rychlost, fine-tuning schopnosti)
% - Tabulka: model × AP × FPS × parametry

\section{Ragdoll constraints}
\label{sec:ragdoll}
% TODO: Post-processing vrstva po pose estimation
% - Limity rotace kloubů (loket 0-150°, koleno 0-170°, atd.)
% - Filtrování fyzikálně nemožných póz (artefakty z okluzí, reflexů)
% - Temporální vyhlazení (výkyvy mezi framy)
% - Detekce záměny levá/pravá strana
% - Důležitost pro podvodní záběry

\section{Biomechanické metriky}
\label{sec:metriky}
% TODO: Úhly kloubů, rotace trupu, frekvence/délka záběru, fáze záběru
% - Tabulka: metrika × úroveň × optimální rozsah × práh chyby
% - 3 úrovně: začátečník / mírně pokročilý / expert

\section{Referenční systém}
\label{sec:referencni_system}
% TODO: 3 úrovně × pohledy × styly
% - Generování referenčních vzorů ze SwimXYZ
% - DTW porovnání (vysvětlení algoritmu)

\section{Porovnávací metoda}
\label{sec:porovnavaci_metoda}
% TODO: DTW + pravidlový systém + skóre 0-100
% - Detekce cyklů záběru
% - Výpočet odchylek od reference
% - Mapování odchylek na konkrétní chyby
% - Generování feedbacku v češtině

\section{Diskuse alternativních přístupů}
\label{sec:diskuse_alternativ}
% TODO: RAG/LLM přístup — proč jsme se rozhodli PROTI
% - Comendant 2024 (RAG coaching for swimming)
% - Talking Tennis (biomechanics → LLM)
% - BoxingPro (IoT → LLM)
% Důvody PROTI:
% - Závislost na externích API (single point of failure)
% - Nedeterministické výstupy (stejný vstup → různý feedback)
% - Latence a náklady na inference
% - Nemožnost plně fungovat offline
% - Obtížná reprodukovatelnost experimentů
% Trade-offs: vyšší flexibilita LLM vs. konzistence a kontrola pravidlového systému

%---------------------------------------------------------------
\chapter{Implementace}
\label{ch:implementace}
%---------------------------------------------------------------

Tato kapitola popisuje implementaci jednotlivých komponent systému SwimAth podle návrhu z~kapitoly~\ref{ch:navrh}.
Kapitola je strukturována od~předzpracování videa přes ML pipeline a~webovou aplikaci až~po~testovací strategii a~nasazení.

\section{Předzpracování videa}
\label{sec:predzpracovani}

% TODO: OpenCV pipeline
% - VideoCapture wrapper: otevření videa, kontrola codec/fps/resolution
% - Validace formátu: MP4/MOV, H.264/H.265 codec
% - Normalizace: resize na max 1080p (zachování aspect ratio), normalizace FPS na 30
% - FFmpeg integrace: konverze nestandardních formátů
% - Frame extraction: sekvenční čtení framů do numpy array
% - Buffering strategie: batch processing po N framech (memory efficient)

\section{ML pipeline}
\label{sec:ml_pipeline}

% TODO: Implementace ML pipeline jako SW komponenty

\subsection{ONNX Runtime inference}

% TODO: Popis inference pipeline
% - Načtení modelu: ONNX Runtime session s GPU ExecutionProvider
% - Batched inference: zpracování N framů najednou (optimalizace GPU utilization)
% - Pre-processing: resize, normalizace, top-down crop podle bounding boxu
% - Post-processing: heatmap → keypoints (argmax + subpixel refinement)
% - Fallback: pokud GPU není dostupný, CPU inference (pomalejší, ale funkční)

\subsection{Strategy pattern pro PE modely}

% TODO: Popis implementace Strategy pattern
% - PoseEstimatorBase (ABC): predict(frames) → keypoints
% - MediaPipePoseEstimator: wrapper pro MediaPipe (baseline)
% - RTMPosePoseEstimator: ONNX inference RTMPose-m
% - ViTPosePoseEstimator: ONNX inference ViTPose-B
% - Výběr strategie na základě konfigurace (env variable / admin setting)
% - Konfigurační rozhraní: model_path, input_size, keypoint_count, confidence_threshold

\subsection{Pipeline orchestrace}

% TODO: Popis pipeline flow
% - Pipeline třída: orchestrace kroků (PE → ragdoll → metrics → DTW → feedback)
% - Každý krok jako samostatná třída s rozhraním process(input) → output
% - Error handling: selhání na jednom framu nezastaví celou analýzu
% - Progress callback: po každém kroku update progressu přes Redis

\section{Backend}
\label{sec:backend}

% TODO: Implementace FastAPI backendu

\subsection{Struktura projektu}

% TODO: Diagram adresářové struktury (listing)
% swimath-api/
% ├── app/
% │   ├── main.py              # FastAPI app, CORS, lifecycle
% │   ├── config.py            # Pydantic BaseSettings
% │   ├── api/                 # Endpointy (upload, analysis, history, websocket)
% │   ├── ml/                  # ML pipeline (PE, biomechanics, comparator, feedback)
% │   ├── processing/          # Video zpracování (handler, extractor, pipeline)
% │   ├── references/          # Referenční databáze (manager, templates)
% │   ├── db/                  # SQLAlchemy modely, Alembic migrace
% │   └── utils/               # Pomocné funkce

\subsection{Routing a dependency injection}

% TODO: FastAPI specifika
% - APIRouter pro modularitu endpointů
% - Dependency injection: DB session, Celery app, config
% - Pydantic modely pro validaci request/response
% - Middleware: CORS, request logging, error handling

\subsection{Pydantic modely}

% TODO: Ukázkové modely (listing)
% - UploadParams: skill_level (Enum), camera_view (Enum), swim_style (Enum)
% - AnalysisResponse: id, status, metrics, deviations, feedback, overall_score
% - ProgressMessage: progress (int), stage (str), message (str)

\section{Frontend}
\label{sec:frontend}

% TODO: Implementace Vue.js 3 frontendu

\subsection{Vue.js 3 Composition API}

% TODO: Architektura frontendu
% - Composition API s <script setup> syntaxí
% - Pinia store pro globální stav (aktuální analýza, historie)
% - Vue Router pro navigaci (upload → progress → results → history)
% - Axios instance s interceptory pro error handling

\subsection{Upload komponenta}

% TODO: Popis upload flow
% - Drag & drop + file input
% - Klientská validace (formát, velikost)
% - Formulář pro výběr parametrů (skill_level, camera_view, swim_style)
% - Axios POST s FormData
% - Redirect na progress view

\subsection{WebSocket progress}

% TODO: Popis real-time progress
% - WebSocket spojení na /ws/progress/{task_id}
% - Reactive progress bar (Vue reactive refs)
% - Stage indikátor (jaký krok pipeline právě běží)
% - Auto-redirect na výsledky po dokončení

\subsection{Video overlay --- Canvas API}

% TODO: Popis vizualizace pózy
% - HTML5 Canvas overlay přes <video> element
% - Vykreslení skeleton (keypoints + connections)
% - Synchronizace canvas s video playback (requestAnimationFrame)
% - Barevné kódování: zelená (OK), žlutá (minor), červená (severe) pro problematické klouby
% - Timeline vizualizace záběrových cyklů

\section{Asynchronní zpracování}
\label{sec:async_impl}

% TODO: Implementace Celery pipeline

\subsection{Celery workers}

% TODO: Konfigurace a implementace
% - Celery app s Redis brokerem
% - Task: analyze_video(analysis_id) — hlavní processing task
% - Konfigurace: concurrency, prefetch multiplier, task time limit
% - GPU worker: dedikovaný worker s GPU přístupem pro ML inference
% - CPU worker: pro lehké úlohy (validace, export)

\subsection{WebSocket progress bridge}

% TODO: Redis pub/sub → WebSocket
% - Worker publikuje progress do Redis kanálu
% - FastAPI WebSocket endpoint odebírá z Redis kanálu
% - Bridge zajišťuje real-time doručení progress updates klientovi

\subsection{Task state management}

% TODO: Stavový diagram tasku
% - Stavy: PENDING, STARTED, PROGRESS, SUCCESS, FAILURE, REVOKED
% - Retry logika: max 3 retries, exponential backoff
% - Timeout: task time limit pro prevenci zombie tasků
% - Cleanup: automatické smazání starých tasků

\section{Databáze}
\label{sec:databaze}

% TODO: Implementace persistence vrstvy

\subsection{SQLAlchemy ORM}

% TODO: Popis ORM modelů
% - Deklarativní mapování (DeclarativeBase)
% - Modely: Analysis, Metric, Deviation, Feedback, ReferenceTemplate
% - Vztahy: relationship() s lazy loading
% - JSON column pro flexibilní metadata

\subsection{Alembic migrace}

% TODO: Migrace workflow
% - alembic revision --autogenerate pro generování migrací
% - alembic upgrade head při startu aplikace
% - Verzování schématu v git

\subsection{Connection pooling}

% TODO: SQLAlchemy connection pool
% - Pool size, max overflow, pool recycle
% - Async session (pokud FastAPI async endpoints)

\section{Referenční databáze}
\label{sec:referencni_db}

% TODO: Implementace referenčního systému

% - Struktura: skill_level / swim_style / camera_view / template.json
% - JSON schéma: temporální sekvence keypointů (DBA výstup) + variance envelope
% - Generování: DBA iterativní algoritmus na trénovacích záběrech ze SwimXYZ
% - Verzování šablon: verze v JSON metadata, migrace při změně formátu
% - Lazy loading: šablony se načítají při prvním požadavku a cacheují v paměti

\section{Testovací strategie}
\label{sec:testovani}

% TODO: Detailní popis testovacího přístupu

\subsection{Testovací pyramida}

% TODO: Diagram testovací pyramidy (figure)
% - Unit testy (základ): jednotlivé funkce, třídy, moduly
% - Integrační testy (střed): API endpointy, DB operace, Celery tasky
% - E2E testy (špička): celý upload → výsledky flow

\subsection{Unit testy}

% TODO: pytest specifika
% - Fixtures pro DB session, test client, mock ML models
% - Parametrizované testy pro různé kombinace (styl × pohled × úroveň)
% - Mock PE modelu pro rychlé testy (předpočítané keypoints)
% - Coverage target: >80%

\subsection{Integrační testy}

% TODO: FastAPI TestClient
% - Test upload endpointu (validace, vytvoření záznamu)
% - Test analysis endpointu (vrácení výsledků)
% - Test WebSocket (progress updates)
% - Test Celery tasku s mock ML modelem

\subsection{E2E testy}

% TODO: Playwright testy
% - Upload flow: nahrání videa → progress → výsledky
% - Video overlay: kontrola renderování
% - Historie: zobrazení předchozích analýz

\section{CI/CD a deployment}
\label{sec:cicd}

% TODO: Popis nasazení a automatizace

\subsection{Docker}

% TODO: Multi-stage Dockerfile
% - Build stage: instalace závislostí, kompilace
% - Runtime stage: minimální image, non-root user
% - docker-compose: frontend, api, worker, db, redis
% - GPU podpora: nvidia-docker pro ML worker

\subsection{CI pipeline}

% TODO: GitHub Actions workflow
% - Lint: ruff, mypy
% - Test: pytest s coverage reportem
% - Build: Docker image
% - (Budoucí) Deploy: push to registry, deploy to server

%---------------------------------------------------------------
\chapter{Experimenty a vyhodnocení}
\label{ch:experimenty}
%---------------------------------------------------------------

\section{Přesnost odhadu pózy na plaveckých datech}
\label{sec:presnost_pe}
% TODO: Benchmark MediaPipe vs RTMPose vs ViTPose na SwimXYZ
% - Metriky: PCK, AP, MPJPE
% - Tabulka výsledků
% - Vliv fine-tuningu

\section{Kvalita detekce chyb}
\label{sec:kvalita_detekce}
% TODO: Porovnání s hodnocením trenéra
% - Metodika evaluace
% - Precision/Recall detekce chyb
% - Confusion matrix

\section{Testování na reálných videích}
\label{sec:realna_videa}
% TODO: Různé úrovně plavců, různé pohledy kamery
% - Autorova vlastní videa (triatlonista)
% - Různé podmínky nahrávání
% - Kvalitativní analýza výstupů

\section{Limity systému}
\label{sec:limity}
% TODO: Kde systém selhává, jaké podmínky potřebuje
% - Minimální kvalita videa
% - Problémy s podvodními záběry
% - Omezení na 4 styly
% - Závislost na kvalitě referenčních vzorů

%---------------------------------------------------------------
\chapter{Závěr}
\label{ch:zaver}
%---------------------------------------------------------------

% TODO: Shrnutí, přínos, budoucí práce
% - Co bylo dosaženo
% - Hlavní přínosy práce
% - Budoucí práce: další styly, více pohledů kamery, mobilní verze


\appendix\appendixinit

%---------------------------------------------------------------
\chapter{Obsah přiloženého média}
\label{app:medium}
%---------------------------------------------------------------

% TODO: Popis obsahu přiloženého média (CD/USB/repozitář)
% - Zdrojové kódy
% - Natrénované modely
% - Ukázková videa
% - Návod na spuštění


\backmatter

\printbibliography[heading=bibintoc]

\end{document}
